\documentclass{article}
\usepackage[T1]{fontenc}
\usepackage{lmodern}
\usepackage{graphicx}
\usepackage{amsmath,amssymb}
\usepackage[a4paper,margin=2cm]{geometry}
\usepackage[absolute,overlay]{textpos}
\setlength{\TPHorizModule}{1mm}
\setlength{\TPVertModule}{1mm}
\usepackage[indonesian]{babel}
\usepackage{hyperref}
\setlength{\emergencystretch}{2em}
% unit: Halaman 13

\begin{document}

% === Halaman 13 (llm) ===

\textbf{Contoh:} $\mathrm{F}_{23}$ mempunyai anggota $\{0, 1, 2, ..., 22\}$.
Contoh operasi aritmetika di dalam $\mathrm{F}_{23}$:
\[ 12 + 20 = 9 \quad (\text{karena } 12 + 20 = 32 \pmod{23} = 9) \]
\[ 8 \cdot 9 = 3 \quad (\text{karena } 8 \times 9 = 72 \pmod{23} = 3) \]

\end{document}
